\documentclass[10pt,a4paper]{amsart}
\usepackage[margin=19mm]{geometry}
\usepackage{amsmath,amssymb,mathtools}
\usepackage[scr=rsfs,cal=euler]{mathalfa}
\usepackage{xfrac}
\usepackage[unicode,hidelinks,bookmarks=false]{hyperref}
\newcommand{\ie}{i.e.\ }
\newcommand{\cf}{cf.\ }
\newcommand{\resp}{respectively\ }
\newcommand{\Subsection}{\subsection}
\providecommand{\nobreakdash}{\nobreak\hbox{-}}
\setlength{\emergencystretch}{2em}
\multlinegap0pt
\usepackage{amssymb}
\usepackage{mathtools}
\usepackage[shortlabels]{enumitem}
\usepackage{xcolor}
\newtheorem{lemma}{Lemma}
\newcommand{\Law}{\operatorname{Law}}
\newcommand{\Prob}{\operatorname{Prob}}
\newcommand{\RR}{\mathbb R}
\title{Log-H\"oder continuity at zero Lyapunov gap for finite-state Markov $GL(2)$-cocycles}
\author{El Hadji Yaya Tall}
\date{Source: arXiv:2608.22157v1 (CC BY 4.0)}
\begin{document}
\maketitle

\noindent\emph{Source and licence.} Log-H\"oder continuity at zero Lyapunov gap for finite-state Markov $GL(2)$-cocycles. Version arXiv:2608.22157v1 (CC BY 4.0), distributed by arXiv under the Creative Commons Attribution 4.0 International licence, http://creativecommons.org/licenses/by/4.0/, as recorded in the arXiv licence metadata for this exact version and shown on the abstract page https://arxiv.org/abs/2608.22157v1. Full licence notice: \texttt{licenses/arxiv-2608.22157-CC-BY-4.0.txt}.\par\smallskip

\noindent\textit{Editorial scope.} Counted unit: the article's lemma lem:path-laws (source0.tex lines 573--584). The article's complete original proof of it is delivered with it (source0.tex lines 586--777, one \texttt{proof} environment of 192 lines). No mathematics is written, repaired or re-derived: the statement and the proof are the article's own lines, in the article's own notation, and no retained line is substituted or re-flowed.  No further source-local context is needed: every cross-reference of every retained line resolves inside the two delivered spans.  Nothing was omitted from any retained span, so the package carries no omission record.  No retained line contains a citation, so the wrapper carries no bibliography and no bibliography entry.  No retained line contains a figure environment, an image-inclusion command or drawing code, so no image is delivered and the package declares no asset dependency.  Editorial formatting only, in addition: an AMS \texttt{amsart} preamble that restates the article's own declarations and macros used by the retained lines, namely the environments \texttt{lemma} on separate counters, together with the standard packages those lines need.  The retained environments print in this document as Lemma 1 in order of appearance, whereas the article prints the same items with the numbering of its own full document.  The complete native source of the article as distributed in the arXiv e-print for this exact version is bundled unchanged as \texttt{sources/208307/source0.tex}.
\begin{lemma}
\label{lem:path-laws}
In a neighborhood of a primitive $P$, the stationary vector
$p(Q)$ depends Lipschitz continuously on $Q$.  Moreover, the
stationary laws of paths $(X_0,\ldots,X_n)$ for $P$ and $Q$ may
be coupled so that
\begin{equation}\label{eq:path-tv}
 d_{\rm TV}\bigl(\Law_P^{[0,n]}(X_0,\ldots,X_n),
                   \Law_Q^{[0,n]}(X_0,\ldots,X_n)\bigr)
 \le C(n+1)\|P-Q\|_1.
\end{equation}
\end{lemma}

\begin{proof}
Write row vectors on the left and let
$$
 H_0=\left\{v\in\RR^N:\sum_iv_i=0\right\}.
$$
Since $P$ is irreducible, the eigenvalue $1$ is simple.  Therefore
the map
$$
 v\longmapsto v(I-P)
$$
is invertible from \(H_0\) onto itself.  If $q=q(Q)$, then
$q(I-Q)=0$, while $p(I-P)=0$.  Subtracting gives
\begin{equation}\label{eq:stationary-vector-linearization}
 (q-p)(I-P)=q(Q-P).
\end{equation}
Both sides belong to $H_0$.  Applying the inverse of $I-P$ on
$H_0$, and using $\|q\|_1=1$, yields
\begin{equation}\label{eq:stationary-vector-lipschitz-expanded}
 \|q-p\|_1\le C_P\|Q-P\|_1.
\end{equation}
The same inverse remains uniformly bounded in a sufficiently small
neighborhood of $P$, proving the first assertion.


Now define
$$
 E_k=\{X_0=Y_0,\ldots,X_k=Y_k\}.
$$
Thus, $E_k$ is the event that the two paths coincide through time
$k$.

We first maximally couple the initial stationary distributions
$$
 X_0\sim p(P),
 \qquad
 Y_0\sim p(Q).
$$
By the defining property of a maximal coupling,
\begin{equation}\label{eq:initial-maximal-coupling}
 \Prob(X_0\ne Y_0)
 =
 d_{\rm TV}\bigl(p(P),p(Q)\bigr)
 =
 \frac12\|p(P)-p(Q)\|_1.
\end{equation}

Now suppose that $E_k$ occurs and that the common state at time $k$ is $i$.  Conditionally on this event, the laws of the next states are
$$
 \Law_P(X_{k+1}\mid E_k,X_k=Y_k=i)=P_{i\bullet},
$$
and
$$
 \Law_Q(Y_{k+1}\mid E_k,X_k=Y_k=i)=Q_{i\bullet}.
$$
Choose a maximal coupling of these two probability vectors.  Then
\begin{align}
 &\Prob\bigl(
   X_{k+1}\ne Y_{k+1}
   \mid E_k,X_k=Y_k=i
 \bigr)
 \notag\\
 &\qquad
 =
 d_{\rm TV}(P_{i\bullet},Q_{i\bullet})
 =
 \frac12\sum_j|P_{ij}-Q_{ij}|.
 \label{eq:one-step-maximal-coupling}
\end{align}
Since
$$
 \|P-Q\|_1
 =
 \max_i\sum_j|P_{ij}-Q_{ij}|,
$$
equation \eqref{eq:one-step-maximal-coupling} gives, uniformly in
\(i\),
\begin{equation}\label{eq:one-step-disagreement-bound}
 \Prob\bigl(
   X_{k+1}\ne Y_{k+1}
   \mid E_k,X_k=Y_k=i
 \bigr)
 \le \frac12\|P-Q\|_1.
\end{equation}

We now estimate the probability that the two complete paths differ.
Define the first disagreement events
$$
 D_{-1}=\{X_0\ne Y_0\},
$$
and, for $0\le k\le n-1$,
$$
 D_k
 =
 E_k\cap\{X_{k+1}\ne Y_{k+1}\}.
$$
Thus, $D_k$ is the event that the paths agree through time $k$ and
separate at the transition from time $k$ to time $k+1$.  These
events are pairwise disjoint, and
\begin{equation}\label{eq:path-disagreement-decomposition}
 \bigl\{
   (X_0,\ldots,X_n)\ne(Y_0,\ldots,Y_n)
 \bigr\}
 =
 D_{-1}\cup\bigcup_{k=0}^{n-1}D_k.
\end{equation}

Using conditional probability and
\eqref{eq:one-step-disagreement-bound}, we obtain
\begin{align}
 \Prob(D_k)
 &=
 \sum_i
 \Prob(E_k,X_k=Y_k=i)
 \notag\\
 &\qquad\quad\times
 \Prob\bigl(
   X_{k+1}\ne Y_{k+1}
   \mid E_k,X_k=Y_k=i
 \bigr)
 \notag\\
 &\le
 \frac12\|P-Q\|_1
 \sum_i\Prob(E_k,X_k=Y_k=i)
 \notag\\
 &=
 \frac12\|P-Q\|_1\Prob(E_k)
 \notag\\
 &\le
 \frac12\|P-Q\|_1.
 \label{eq:first-disagreement-bound}
\end{align}

Therefore, by \eqref{eq:initial-maximal-coupling},
\eqref{eq:path-disagreement-decomposition}, and
\eqref{eq:first-disagreement-bound},
\begin{align}
 &\Prob\bigl(
   (X_0,\ldots,X_n)\ne(Y_0,\ldots,Y_n)
 \bigr)
 \notag\\
 &\qquad
 =
 \Prob(D_{-1})
 +\sum_{k=0}^{n-1}\Prob(D_k)
 \notag\\
 &\qquad
 \le
 \frac12\|p(P)-p(Q)\|_1
 +\frac n2\|P-Q\|_1.
 \label{eq:path-disagreement-expanded}
\end{align}
Using the Lipschitz estimate for the stationary distributions,
$$
 \|p(P)-p(Q)\|_1
 \le C_P\|P-Q\|_1,
$$
we conclude that
\begin{align}
 \Prob\bigl(
   (X_0,\ldots,X_n)\ne(Y_0,\ldots,Y_n)
 \bigr)
 &\le
 \frac{C_P+n}{2}\|P-Q\|_1
 \notag\\
 &\le
 C(n+1)\|P-Q\|_1.
 \label{eq:path-coupling-expanded}
\end{align}

After the two chains separate, they may be coupled in an arbitrary
way. %  Indeed, once they differ at some coordinate, the event on the
%left-hand side of \eqref{eq:path-coupling-expanded} has already
%occurred; their subsequent behavior cannot make the two complete path
%vectors equal.
%
Finally, since the joint law of
$$
 \bigl((X_0,\ldots,X_n),(Y_0,\ldots,Y_n)\bigr)
$$
is a coupling of the two path laws, the coupling inequality gives
$$
 d_{\rm TV}\bigl(
   \Law_P^{[0,n]}(X_0,\ldots,X_n),
   \Law_Q^{[0,n]}(Y_0,\ldots,Y_n)
 \bigr)
 \le
 C(n+1)\|P-Q\|_1.
$$



\end{proof}

\end{document}
